\subsection{An actual integral resolution and integral transfer of the quadratic repair}

\begin{theorem}[AT22: Proved; integral repair of the specified quadratic secondary class]
\label{AT22}
Let $X\to A,B$ be actual strict maps over $S_{\mathbb Z}$,
with torsionfree terms, $X_0=S$, $H^0(X_n)=S$, and rational
coefficient comparison in every simplex and weight.
Let $A,B$ be admissible.
Suppose first that normalized marks satisfy
$\partial h=\rho\in S^b$ and $\partial x_0=dh$ as in
(AT.116), with $b\geq3$ and $s=2b$.
Choose $N=2^e$ satisfying the first two bounds of (AT.117),
and suppose in addition
\[
 \begin{split}
 s+1&\geq
  \max_{C=A,B}\{\lambda_C(4),\mu_C(4)\},\\
 s&\geq
  \max_{C=A,B}\{\lambda_C(2),\mu_C(2)\}.
 \end{split}                                                     \tag{AT.150}
\]
There is an explicitly constructed integral graded-free
coefficient resolution containing the displayed low generators
of A-S09-22. Use it in AT16 with
$L=I_1\langle s\rangle$, $m=2$,
$z_0=\gamma_2(dx_0)$ and $\eta_0=x_0dx_0$.
For every resulting pair of actual target lifts, the specific
corrected class $\zeta_{\rm tw}$ has zero integral relative
period modulo two and admits the actual simultaneous
integral repair (AT.107).
The same repair extends the actual target maps after
A-SAT01 saturation.

For an original C21 cell with arbitrary one-form boundary
over the free even-disk first skeleton, use the explicit
correction of AT21 and exactly the original edge choice
(AT.140). With the same bounds and all the same source
and target hypotheses, the original resolved and degree-three
repair algebras have the corresponding actual integral
maps. This conclusion uses explicit integral inverse maps;
no local minimal resolution is asserted to descend.

The theorem concerns the displayed quadratic secondary
class and its repair. It proves neither a nonzero integral
cohomology class for the newly chosen resolution nor an
eventual vertical comparison after all repairs.
\end{theorem}

\begin{proof}
We first construct the resolution over $\mathbb Z$ itself.
Before the shift, take
\[
 Q_0=\bigoplus_{a\geq1}S e_a,\qquad
       \operatorname{wt}(e_a)=a,\qquad
       \varepsilon(e_a)=T^{[a]}.
 \tag{AT.151}
\]
This is surjective onto $I_1$ and has finite free groups
in every weight.
Put $U=T^{[2]}$.
In $Q_1$ include distinct free generators $r_2,r_3$
with the actual cycle boundaries
\[
                 \partial r_2=Te_1-2e_2,\qquad
                 \partial r_3=Ue_1-Te_2.
 \tag{AT.152}
\]
The augmentation of the second is zero because both
products are $3T^{[3]}$.
In $Q_2$ include distinct free generators $g_3,g_4$
and in $Q_3$ include $h_4$, with
\[
 \partial g_3=Tr_2-2r_3,\qquad
 \partial g_4=Ur_2-Tr_3,\qquad
 \partial h_4=Tg_3-2g_4.
 \tag{AT.153}
\]
Every displayed boundary is a cycle over $\mathbb Z$;
this follows by direct substitution and $T^2=2U$.

Complete these choices as follows. In each homological
degree and each weight, the previous map has a kernel
which is a subgroup of a finite free abelian group.
It is therefore finitely generated and free.
Choose finitely many homogeneous generators spanning
the quotient of that kernel by the boundaries already
supplied from lower weights. Such a finite generating
set exists even when that quotient has torsion.
Adjoin their free $S$ generators, retaining each of the
specified generators (AT.152)--(AT.153), whether or not
it is redundant. At its specified weight add those
named cycles before choosing the remaining generators.
Induction on weight makes the map onto the entire
kernel; induction on homological degree constructs
a free resolution. Only finitely many generators
occur in a fixed weight in each degree, so every
weight group of every term and kernel stays finite
free. All generators have positive weight.
Now apply the common shift $s$.

This is an actual integral resolution containing the
literal named chain. It is not the assertion that a
chosen minimal local resolution lifts.
Its cone and truncation $P,W$ are exactly those of
AT16. Define
\[
 C=[v_{g_3}],\quad H=[u_{h_4}],\quad K=[v_{g_4}],
 \quad g=u_{g_3},\quad l=u_{g_4}.
\]
The boundary of $v_{h_4}$ gives $TC=2(K-H)$.
The closed element $-Ug+Tl$ of $W^1$ defines the actual
$J\in W^0=\ker(W^1\to W^2)$ with that derivative.
Injectivity of $d:W^0\to W^1$ proves $TH=2J$.
The same torsionfree calculation gives $TK-UC=2J$.
Thus every displayed polynomial defining $B$ and
$\zeta_{\rm tw}$ is present in this resolution and
satisfies $dB=2\zeta_{\rm tw}$ by the proved generator
identities. Extra resolution generators do not enter
or alter those identities.

The first bound in (AT.150) permits the actual integral
AT16 construction and its target maps using this free
resolution. Fix any such lift to $A$.
The natural period comparison in weight $2s+5$ identifies
its residue with
\[
          \beta_A=[f_A(B)\bmod2]\in S^{2s+5}/2S^{2s+5}.
 \tag{AT.154}
\]
We prove that this actual residue is zero.

Localize all displayed modules and maps at two.
Exactness is preserved because $\mathbb Z_{(2)}$ is flat.
We therefore have the particular free resolution
$Q\otimes\mathbb Z_{(2)}$ with its original augmentation
and all its named generators.
The target bounds remain valid after localization.
Repeat the construction of AT20 on this resolution:
its local coefficient model and the elements $q,\alpha$
construct the normalized quadratic $D(3)$ data;
every free generator still has weight at least $s+1$;
the total-degree-four Hom vanishing and the terminal
equation therefore give actual resolved generator values
on this same $W\otimes\mathbb Z_{(2)}$.
The correction (AT.129), followed by (AT.131), supplies
a local target lift with zero residue.
This argument uses only the free resolution and its
augmentation, not its minimality or any missing
comparison between two resolutions.

We verify explicitly that the choice-independence
calculation applies to the present resolution too.
Differences of its target lifts are classes in
\[
 \operatorname{Ext}^1_{\mathrm{gr}S_{(2)}}
       ((I_1\langle s\rangle)/2,S_{(2)})
                   =S_{(2)}^s/2S_{(2)}^s.
 \tag{AT.155}
\]
The actual free Hom comparison is the finite diagonal
of AT16, and the last equality is the long exact
sequence for multiplication by two on the actual
coefficient tail, using AT04.
No comparison with a minimal resolution is used.
The second bound in (AT.150) supplies a closed relative
two-form of weight $s$ and coefficient period $T^{[s]}$.
Its coefficient multiples on $v_0$, zero on all other
summands, represent every class in (AT.155).
They change $B$ by $T\omega H=2J\omega$.
After subtracting such a representative, the difference
is an actual Hom boundary. Formula (AT.99) still applies:
it uses only the named differential identities
(AT.152)--(AT.153) and $TC=2(K-H)$, not a list of all
resolution generators. It changes $B$ modulo two by
an actual relative boundary.
Thus every lift on this particular localized resolution,
including the localization of the chosen integral one,
has the same zero residue.

Finally this vanishing descends by an actual comparison
of complexes. The relative module
$\mathcal C_A=\operatorname{Hom}(K(2),A)$ is the kernel
of the three face maps from $A_2$ to $A_1^3$.
Flat localization commutes with that finite kernel, so
\[
 \mathcal C_A\otimes\mathbb Z_{(2)}
           \xrightarrow{\;\cong\;}\mathcal C_{A_{(2)}}.
\]
Reduction modulo two gives an isomorphism already
before taking cohomology,
\[
               \mathcal C_A/2\mathcal C_A
                 \xrightarrow{\;\cong\;}
               \mathcal C_{A_{(2)}}/2\mathcal C_{A_{(2)}} .
 \tag{AT.156}
\]
Indeed every odd integer acts invertibly on every
module killed by two. The coefficient-period
comparisons commute with these maps by their
displayed natural exact-sequence construction.
Also $S^w/2\to S_{(2)}^w/2$ is an isomorphism.
Consequently (AT.154) is zero integrally.

For completeness the primitive construction works
over $\mathbb Z$ itself. The period identification
identifies (AT.154) with the zero cohomology class
of $f_A(B)$ in $\mathcal C_A/2\mathcal C_A$.
Choose an integral $v\in\mathcal C_A^1$ whose
differential is congruent to $f_A(B)$ modulo two, and set
\[
        U_A=\frac{f_A(B)-dv}{2}\in\mathcal C_A^2,
        \qquad V_A=-v .
\]
The first expression is integral by the congruence.
The identity $dB=2\zeta_{\rm tw}$ and torsionfreeness
give $dU_A=f_A(\zeta_{\rm tw})$, while
$dV_A=2U_A-f_A(B)$ holds directly.
These are exactly (AT.107), proved over $\mathbb Z$
rather than by invoking the locally stated AT18.
The same argument for the independently chosen
integral lift to $B$ gives $U_B,V_B$ on the very
same source and repair. A-SAT01 at cutoff $s+1$
constructs the actual saturation with those target
maps; the already specified values of $B,\zeta_{\rm tw}$
are unchanged, so the same $U,V$ extend its maps.

For the final assertion apply AT21 to the old cell.
The corrected mark $x^\flat=x-t$ has the required
boundary $dh$. Apply the argument just proved to
that actual mark and the integral resolution (AT.151)--(AT.153).
Now the literal named generators and $J$ have been
constructed over $\mathbb Z$, so the calculations
(AT.142)--(AT.144) apply over $\mathbb Z$ to this
particular resolution. The explicit chosen edges
(AT.140), inverse maps (AT.141), and repair translation
(AT.144) give the original-source maps and repair.
Their saturation comparison is the actual
stagewise root transport proved in AT21b.
This proves the stated integral conclusion with all
morphisms specified.
\end{proof}
